% ---------------------------------------------------------------------------
% Abstract (input by main.tex after \maketitle; no preamble).
% Numbers: final rerun of 2026-10-12 (paper/results/*.json): network.json (China leads 10 of 18
% basket stages; median designed HHI 0.1520; share of production outside the three largest
% exporting nodes, median stage 0.6625; non-high-income Upstream share 0.465), strategic.json and
% robustness.json (27 Batteries Midstream picks; 21 of 57 margins below 0.01), method.json.
% ---------------------------------------------------------------------------
\begin{abstract}
\noindent No allied economy can efficiently reproduce the whole supply chain of energy
technologies, so protection against concentrated suppliers must come from allocating
capabilities across trusted partners: managed interdependence rather than autarky. We build a
framework to measure and design it. Energy-security and economic-opportunity pillars
score 219 economies in up to 33 technology--stage cells. A Strategic Index reads security as need,
and an endowment screen admits extraction stages to a member's priority only where it has a
resource base. A 57-member coalition (37 network
nodes) is derived from UN General Assembly ideal points, agreements and world supply shares.
Bilateral indices rate partners as suppliers and markets, and a mixed-integer program allocates
production in 22 upstream and midstream stages under a diversification penalty calibrated to
observed export concentration. China is the largest exporter in 10 of the 18 traded stages. The
calibrated design spreads each stage over seven to ten producers (median Herfindahl index 0.152)
and, because it scores capability rather than current capacity, places two-thirds of the median
stage's production outside today's three largest exporting nodes and 47 percent of upstream output
in non-high-income members. Battery midstream is the most common national priority (27 members),
far more than the nine producers the network design assigns it, and 21 of the 57 priorities lie
within 0.01 of the runner-up. The
framework turns rankings into an auditable division of labour.
\end{abstract}

\medskip
\noindent\textbf{JEL codes:} C43, C61, F14, F52, L52, Q48.

\noindent\textbf{Keywords:} allied industrial base; energy security; economic opportunity;
composite indicators; friend-shoring; geopolitical alignment; supply-chain network design;
industrial policy.
